% Frozen snapshot: <writing-loop-fixture, 2>
\documentclass{article}
\usepackage{amsmath}
\begin{document}

\begin{abstract}
Reliable manuscript construction requires more than grammatical correction.
This fixture represents a compact journal structure in which the chapter
sequence, local argument, notation, equations, terminology, and schematic
handoff are checked after every substantive edit. The method uses a frozen
conception context to keep each section tied to one scientific question and
one reader-facing result, while its notation registry defines every
mathematical object before use. A subsection audit then records the logical
role of adjacent sentences, the causal chain from the motivating limitation to
the proposed mechanism and its consequence, and the assumptions and dimensions
needed by each formula. It also verifies that every retained sentence has an
indispensable purpose and that required principle, model-structure, and
task-workflow schematics have a consistent LaTeX handoff. The resulting
workflow provides a repeatable record of revisions and blocks delivery whenever
an unresolved conflict remains. This record also supports transparent review
before figure and experiment work.
\end{abstract}

\section{Introduction}

The main contributions are:

\begin{enumerate}
\item A compact chapter-control procedure aligns each central problem with one
principal contribution and one explicit body result.
\item A ten-loop subsection audit and paragraph-purpose map record structural,
logical, linguistic, symbolic, mathematical, and schematic evidence after
every substantive edit.
\end{enumerate}

\section{Preliminaries and Problem Formulation}

\subsection{Problem Formulation}

\begin{enumerate}
\item Maintain a scientifically ordered chapter sequence without duplicated
section responsibilities.
\item Verify every revised subsection through the complete writing-quality
loop before advancing.
\end{enumerate}

\section{Proposed Method}

Let $x$ denote the scalar model state and let $y$ denote the scalar output. The
complete fixture model is

\begin{equation}
  y=x.
  \label{eq:writing-loop-fixture}
\end{equation}

Because the frozen baseline determines the scientific order before drafting,
the ten-loop audit can examine each substantive revision without silently
changing the paper's task. Its LaTeX schematic record then supplies the
topological handoff for Stage 4 while making no experimental claim.

\section{Conclusion}

The chapter and subsection records provide the evidence required by the
writing loop.

\end{document}
